%% ============================================================
%%  CHAPTER 7 — PACK-LEVEL SIMULATION
%% ============================================================
\chapter{Pack-Level Simulation}
\label{ch:pack}

\section{Pack Topology and Configuration}

The \VIBE{} framework supports arbitrary $N_s \times N_p$ rectangular pack topologies, where $N_s$ is the number of cells in series and $N_p$ is the number of parallel branches. Cell $(s,p)$ occupies position $(s, p)$ in the grid, with a linear index $k = s \cdot N_p + p$. Series groups share current; parallel branches share voltage.

\section{Heterogeneous Cell Parameters}
\label{sec:hetero_params}

Each cell can carry independent physical parameters. The \texttt{overrides} dictionary in \texttt{ImplicitBatterySolver} allows per-cell specification of any parameter:

\begin{lstlisting}[language=Python, caption={Heterogeneous parameter overrides}]
overrides = {
    (s_idx, p_idx): {
        "hA": WEAK_COOLING_HA,       # 10x lower cooling for branch p
        "R_contact": R_CONTACT_20PCT, # 20% contact resistance mismatch
    }
    for s_idx in range(N_SERIES)
}
\end{lstlisting}

Parameters are stored as tensors of shape $[N_\text{cells}, 1]$ in \texttt{BatteryPhysics.params}, so the vectorised computations automatically use cell-specific values.

\section{Current Distribution Algorithm}
\label{sec:current_dist_algorithm}

The current distribution solver is called at the beginning of each Newton timestep to compute the per-cell current $I_k$ given the pack current $I_\text{pack}$.

\subsection{Algorithm}

For each series row $s$, the following iteration is performed:

\begin{algorithm}[H]
\caption{Kirchhoff Current Distribution with Butler--Volmer Dynamic Resistance}
\label{alg:current_dist}
\begin{algorithmic}[1]
\State Compute circuit parameters: $\text{OCV}, R_\text{series}, I_{0,n}, I_{0,p}, T, V_\text{conc}$ via \texttt{get\_circuit\_parameters()}
\State Initialise: $I_{s,p} \leftarrow I_\text{pack}/N_p$ for all $p$
\For{iter $= 1$ to $5$}
  \State Compute dynamic resistance: $R_{d,n} = 2\mathcal{R}T/\mathcal{F} / \sqrt{I^2 + 4I_{0,n}^2}$
  \State Compute dynamic resistance: $R_{d,p} = 2\mathcal{R}T/\mathcal{F} / \sqrt{I^2 + 4I_{0,p}^2}$
  \State $R_\text{tot} \leftarrow R_\text{series} + R_{d,n} + R_{d,p}$
  \State $V_\text{virt} \leftarrow \text{OCV} - I \cdot R_\text{series} - 2\frac{\mathcal{R}T}{\mathcal{F}}\sinh^{-1}\!\frac{I}{2I_{0,n}} - 2\frac{\mathcal{R}T}{\mathcal{F}}\sinh^{-1}\!\frac{I}{2I_{0,p}} - V_\text{conc}$
  \State $G \leftarrow 1/R_\text{tot}$
  \State $V_\text{row} \leftarrow \left(\sum_p V_\text{virt,p} G_p - I_\text{pack}\right) / \sum_p G_p$
  \State $I_{s,p} \leftarrow (V_\text{virt,p} - V_\text{row}) G_p$
\EndFor
\State \Return $\mathbf{I} \in \mathbb{R}^{N_\text{cells}}$
\end{algorithmic}
\end{algorithm}

\subsection{Convergence and Accuracy}

The conductance-weighted iteration converges quadratically in the vicinity of the solution because it is equivalent to a linearised Newton step on the Kirchhoff equations. Five iterations are sufficient for $<0.1\%$ current error at typical operating conditions.

\section{Pack Voltage Calculation}
\label{sec:pack_voltage}

The pack terminal voltage is:
\begin{equation}
  V_\text{pack} = \sum_{s=1}^{N_s} \left[\frac{1}{N_p}\sum_{p=1}^{N_p} V_{\text{cell},s,p}(I_{s,p})\right]
  \label{eq:pack_voltage}
\end{equation}
i.e., the sum of the mean cell voltages in each series group. This is the standard definition for a rectangular pack topology.

\section{Thermal Interconnection Matrix Construction}
\label{sec:thermal_matrix_construction}

The conductance matrix $\boldsymbol{G} \in \mathbb{R}^{N_\text{cells} \times N_\text{cells}}$ is assembled by the \texttt{\_build\_thermal\_connection\_matrix()} method:

\begin{lstlisting}[language=Python, caption={Thermal connection matrix assembly}]
def _build_thermal_connection_matrix(self, k_contact=0.5, area=0.005, dist=0.02):
    G = torch.zeros((self.n_cells, self.n_cells))
    G_val = k_contact * area / dist          # [W/K] = 0.125 W/K
    for s in range(self.n_series):
        for p in range(self.n_parallel):
            idx = s * self.n_parallel + p
            # Parallel neighbours (same series group)
            if p < self.n_parallel - 1:
                G[idx, idx+1] = G[idx+1, idx] = G_val
            # Series neighbours (adjacent groups)
            if s < self.n_series - 1:
                G[idx, idx+self.n_parallel] = G[idx+self.n_parallel, idx] = G_val
    return G
\end{lstlisting}

The resulting matrix is symmetric, tridiagonal (in the pack index ordering), and has row sums of $\sum_j G_{kj} = n_\text{adj} \cdot G_\text{val}$ where $n_\text{adj} \in \{1,2,3,4\}$ is the number of adjacent cells (corner, edge, or interior).

\subsection{Physical Interpretation}

The conduction term in the thermal ODE (\cref{eq:Q_cond}) can be written in matrix form as:
\begin{equation}
  \mathbf{Q}_\text{cond} = (\boldsymbol{G} - \text{diag}(\boldsymbol{G}\mathbf{1})) \boldsymbol{T} = \boldsymbol{L}_G \boldsymbol{T}
  \label{eq:laplacian_G}
\end{equation}
where $\boldsymbol{L}_G$ is the graph Laplacian of the pack thermal network. The eigenvalues of $\boldsymbol{L}_G$ control the thermal equilibration timescales; the smallest nonzero eigenvalue gives the slowest thermal mode.

\section{Heterogeneity Study: 5S5P Pack}
\label{sec:5s5p}

The \texttt{test\_5s5p\_heterogeneity.py} script implements a literature-backed heterogeneity scenario~\cite{Gogoana2014}:

\begin{table}[H]
\centering
\caption{Heterogeneity parameters for the 5S5P pack study.}
\label{tab:hetero_params}
\begin{tabular}{lcc}
\toprule
\textbf{Parameter} & \textbf{Nominal branches (p=0,1,3,4)} & \textbf{Heterogeneous branch (p=2)} \\
\midrule
$hA$ [\si{\watt\per\kelvin}] & 0.531 (10$\times$ enhanced) & 0.0531 (baseline) \\
$R_\text{contact}$ [\si{\ohm}] & 0 & $0.2\times R_\text{ohm,0} = 1.14\times10^{-3}$ \\
Weakest cell $(s=2,p=2)$ & -- & $R_\text{contact} = 0.3 \times R_\text{ohm,0}$ \\
\bottomrule
\end{tabular}
\end{table}

This combination of reduced cooling ($10\times$ lower $hA$) and elevated contact resistance ($20$--$30\%$ mismatch) in branch $p=2$ creates a realistic heterogeneous scenario. The thermal--electrochemical coupling propagates this heterogeneity: the hotter, higher-resistance branch draws less current, redistributing current to neighbouring branches, which in turn run cooler but carry higher current---a complex multi-physics interaction not capturable by ECM models.

%%------------------------------------------------------------
\section{Comparison with liionpack: MNA vs.\ Nonlinear KCL}
\label{sec:liionpack_comparison}
%%------------------------------------------------------------

\subsection{The liionpack Architecture}
\label{ssec:liionpack_arch}

liionpack~\cite{Tranter2022} is the most prominent open-source pack simulation framework coupling a physics-based cell model (via PyBaMM~\cite{Sulzer2021}) with an electrical network solver. Rather than an iterative KCL loop, it represents the pack topology as a \emph{netlist} of resistors and voltage sources and solves the resulting circuit exactly using \textbf{Modified Nodal Analysis (MNA)}~\cite{Ho1975}---the same sparse-linear formulation used in SPICE.

In MNA, every circuit element is stamped into a sparse conductance matrix $\mathbf{G} \in \mathbb{R}^{n \times n}$, a voltage-source incidence matrix $\mathbf{B} \in \mathbb{R}^{n \times m}$, and source vectors $(\mathbf{i}_s, \mathbf{e}_s)$. The full system is:
\begin{equation}
  \underbrace{\begin{bmatrix} \mathbf{G} & \mathbf{B} \\ \mathbf{B}^\top & \mathbf{0} \end{bmatrix}}_{\mathbf{A}}
  \begin{bmatrix} \mathbf{v} \\ \mathbf{i}_V \end{bmatrix}
  =
  \begin{bmatrix} \mathbf{i}_s \\ \mathbf{e}_s \end{bmatrix}
  \label{eq:mna_system}
\end{equation}
where $n$ is the number of nodes and $m = N_s N_p$ is the number of cells (voltage sources). liionpack assembles $\mathbf{A}$ once per topology and solves $\mathbf{A}\mathbf{x} = \mathbf{z}$ via \texttt{scipy.sparse.linalg.spsolve} at each timestep, updating only $\mathbf{z}$ as the cell voltages evolve.

The liionpack netlist assigns the following element types and default values (taken directly from \texttt{liionpack/netlist\_utils.py}, function \texttt{setup\_circuit}~\cite{Tranter2022}):

\begin{table}[H]
\centering
\caption{liionpack netlist element types and default values; \VIBE{} equivalents are shown for comparison.}
\label{tab:liionpack_resistances}
\begin{tabular}{lllll}
\toprule
\textbf{Descriptor} & \textbf{Physical meaning} & \textbf{Default} & \textbf{Unit} & \textbf{In \VIBE{}} \\
\midrule
$R_b$  & Busbar segment (per gap)   & $10^{-4}$ & $\Omega$ & $R_\text{bus}=10^{-4}$~$\Omega$ \\
$R_c$  & Tab-to-busbar weld         & $10^{-2}$ & $\Omega$ & $R_\text{contact}=10^{-2}$~$\Omega$ \\
$R_t$  & Terminal lug               & $10^{-5}$ & $\Omega$ & absorbed into $R_\text{bus}$ \\
$R_i$  & Cell internal resistance   & from PyBaMM & $\Omega$ & from BV + electrolyte physics \\
\bottomrule
\end{tabular}
\end{table}

\subsection{The Fundamental Difference: Linear vs.\ Nonlinear Cell Model}
\label{ssec:linear_vs_nonlinear}

Each cell appears in the liionpack netlist as a time-varying voltage source:
\begin{equation}
  V_k^\text{LP}(t) = V_{\text{OCV},k}(t^{n-1}) + R_{i,k}(t^{n-1}) \cdot I_k
  \label{eq:liionpack_cell_model}
\end{equation}
where $V_{\text{OCV},k}$ and $R_{i,k}$ are queried from the PyBaMM model at the \emph{previous} timestep $t^{n-1}$ and then held fixed during the MNA solve for step $t^n$. Consequently, the MNA system \cref{eq:mna_system} is linear in all unknowns and yields the current distribution in a single solve. However, the Butler--Volmer exchange current $I_{0,k}(c_{s,k}, c_{e,k}, T_k)$---which depends on the current-step state---does not enter the circuit solve; it is absorbed into the frozen $R_{i,k}$.

In \VIBE{}, the cell model inside Algorithm~\ref{alg:current_dist} is the \emph{full} Butler--Volmer relationship evaluated at the \emph{current-step} state:
\begin{equation}
  V_k^\text{VB}(I_k) = V_{\text{OCV},k}^n
    - I_k R_{\text{series},k}^n
    - \frac{2\mathcal{R}T_k^n}{\mathcal{F}}\sinh^{-1}\!\frac{I_k}{2I_{0,n,k}^n}
    - \frac{2\mathcal{R}T_k^n}{\mathcal{F}}\sinh^{-1}\!\frac{I_k}{2I_{0,p,k}^n}
    - V_{\text{conc},k}^n
  \label{eq:vibe_cell_model}
\end{equation}
The superscript $n$ indicates that all quantities are evaluated at the current Newton state $\mathbf{y}^n$, not at the previous step. This coupling is what makes the circuit nonlinear and requires the iterative solver of \cref{alg:current_dist}.

\subsection{Is the VIBE Iteration Correct?}
\label{ssec:bv_correctness}

The fixed-point iteration is equivalent to a Newton step on the nonlinear Kirchhoff residual:
\begin{equation}
  F_p(I_{s,p}) \equiv V_k^\text{VB}(I_{s,p}) - V_\text{row} = 0,
  \quad \text{subject to} \quad \sum_{p=1}^{N_p} I_{s,p} = I_\text{pack}
  \label{eq:kirchhoff_nonlinear}
\end{equation}
Since $V_k^\text{VB}$ is monotonically decreasing and strictly convex in $I_k$ (the $\sinh^{-1}$ is concave; the sign reversal makes it convex for $I_k > 0$), $F_p$ has a unique zero and the Banach fixed-point theorem guarantees convergence of the iteration when initialised within the basin of attraction (which includes $I_{s,p} = I_\text{pack}/N_p$). In practice, five iterations reduce the KCL residual to below $0.1\%$ of $I_\text{pack}$ for all tested scenarios.

This is \\emph{mathematically exact} for the nonlinear Kirchhoff system---no approximation is introduced---so the \VIBE{} algorithm is not an approximation relative to liionpack; it solves a richer problem.

\subsection{Quantifying the Accuracy Advantage at High C-Rates}
\label{ssec:accuracy_high_c}

The error introduced by liionpack's linearisation \cref{eq:liionpack_cell_model} can be bounded. For a single cell with Butler--Volmer kinetics, the terminal voltage error when using a frozen resistance $R_i$ instead of the true $\sinh^{-1}$ relationship is:
\begin{equation}
  \delta V_\text{BV}(I) = \frac{2\mathcal{R}T}{\mathcal{F}}\sinh^{-1}\!\frac{I}{2I_0}
    - R_{i,\text{frozen}} \cdot I
  \label{eq:bv_error}
\end{equation}
At low C-rates ($I \ll 2I_0$), $\sinh^{-1}(I/2I_0) \approx I/2I_0$ and the linearisation is valid. At 2C--3C, where $I/(2I_0) \sim 1$--$3$, the error $\delta V_\text{BV}$ reaches $O(10\text{--}30)$~mV per cell. Crucially, this voltage error translates directly into a current distribution error: if cell $k$ has a higher-than-actual apparent resistance, the MNA solver will assign it a lower current, and this lower current feeds into the SEI growth rate $\dot{L}_{\text{SEI},k} \propto I_k$ at the next step. The error therefore compounds across cycles, with each cycle's SEI growth error depending on the accumulated current error from previous cycles. Over 200 cycles, this could amount to a systematic underestimation of degradation in the fastest-discharging cell by several percent.

\subsection{Busbar Positional Asymmetry: liionpack's Structural Advantage}
\label{ssec:busbar_asymmetry}

For wide parallel strings ($N_p \gg 1$), liionpack's explicit netlist is structurally superior to \VIBE{}'s lumped resistance. In liionpack, each busbar gap between adjacent tabs is an independent $R_b$ resistor; the total busbar resistance to cell $p$ from the left terminal is:
\begin{equation}
  R_{\text{bus},p}^\text{LP} = p \cdot R_b \qquad (p = 0, \ldots, N_p-1)
  \label{eq:busbar_positional}
\end{equation}
This positional gradient produces a current mismatch even among identical cells. Tranter et al.~\cite{Tranter2022} show that for $N_p = 16$ and $R_b = 10^{-4}$~$\Omega$, the current non-uniformity is up to 20\%.

In \VIBE{}, every cell receives the same $R_\text{bus}$, so positional busbar asymmetry is not captured. This is a recognised limitation. The extension is straightforward: replace the scalar $R_\text{bus}$ with a position-dependent vector $R_{\text{bus},p} = p \cdot R_b$, computed from the pack layout at initialisation. This is identified as a future direction in \cref{ch:conclusions}.

\subsection{Selective Output: Inspired by liionpack}
\label{ssec:selective_output_ch7}

One engineering best-practice adopted from liionpack is the \textbf{selective output system}. liionpack accepts an \texttt{output\_variables} list and computes only those quantities at each timestep, avoiding large intermediate writes to disk. The \VIBE{} framework mirrors this with the \texttt{OutputSpec} class: users declare which of the following variables to save to the compressed \texttt{simulation\_results.npz}:

\begin{itemize}\setlength\itemsep{0pt}
  \item Per-cell: \texttt{TermV}, \texttt{Curr}, \texttt{SOC}, \texttt{Temp}, \texttt{SEI\_Thick}, \texttt{CapFade}
  \item Overpotentials: \texttt{OCV}, \texttt{Rxn}, \texttt{OhmS}, \texttt{OhmE}, \texttt{Conc}, \texttt{SEI}
  \item External circuit: \texttt{OhmRC} ($= I_k R_\text{contact}$), \texttt{OhmRB} ($= I_k R_\text{bus}$)
  \item Pack-level: \texttt{PackVoltage}, \texttt{PackCurrent}
\end{itemize}

The default preset \texttt{OutputSpec("default")} saves all of the above. A minimal preset \texttt{OutputSpec(["terminal\_voltage","temperature"])} produces the smallest possible output file.

\subsection{Framework Summary}
\label{ssec:framework_summary}

\begin{table}[H]
\centering
\caption{Architectural comparison of \VIBE{} and liionpack.}
\label{tab:framework_comparison}
\begin{tabular}{p{3.5cm}p{5.2cm}p{5.2cm}}
\toprule
\textbf{Aspect} & \textbf{liionpack~\cite{Tranter2022}} & \textbf{\VIBE{} (this work)} \\
\midrule
Cell electrochemistry   & PyBaMM SPM / SPMe / DFN (external)
                        & Custom Chebyshev/FVM DFN \\[4pt]
Circuit solver          & MNA sparse-direct (exact for linear circuit, one solve per step)
                        & Nonlinear KCL fixed-point (5 BV evaluations) \\[4pt]
Cell model in circuit   & Linearised $V_k = V_\text{OCV} + R_i I_k$ (frozen from prev.\ step)
                        & Full $\sinh^{-1}$ BV at current-step state \\[4pt]
Exchange current        & Lagged by one timestep
                        & Evaluated inside the distribution loop \\[4pt]
Busbar topology         & Per-segment $R_b$ (positional asymmetry)
                        & Uniform $R_\text{bus}$ per cell (lumped) \\[4pt]
Thermal coupling        & External / post-processed
                        & Implicit DAE state variable \\[4pt]
GPU parallelism         & No
                        & Yes (\texttt{torch.vmap + jacrev}) \\[4pt]
Selective output        & \texttt{output\_variables} list
                        & \texttt{OutputSpec} class \\
\bottomrule
\end{tabular}
\end{table}

The two frameworks are \emph{complementary}. liionpack excels at exact external circuit modelling and pack topology studies; \VIBE{} excels at tightly-coupled electrochemical-electrical integration and long-cycle degradation studies where the compounding of kinetic errors matters most.
